% Modernised reading — fonds Alexandre Grothendieck, shelfmark 90 (pages 1-24,
% the whole folder).
%
% This is an INTERPRETATION, not a transcription. It re-reads the pages in
% today's notation and vocabulary, states the hypotheses the manuscript leaves
% implicit, and drops the critical apparatus so that the mathematics reads
% straight through. Everything the brackets and underlines carried moves into
% a footnote. It is held to being mathematically correct as it stands; where
% the manuscript is loose or mistaken the true statement appears and the
% footnote says what the page has. In any dispute of substance the
% transcription governs: whoever cites Grothendieck cites batch-01.fr.tex
% (pp. 1-20) or batch-02.fr.tex (pp. 21-24).
%
% Pass: Opus 5.5 (claude-opus-5-5), 2026-10-10 — first pass, unchecked against the pages by a human.
% The folder was read whole, its two batches together; this is one document
% for the shelfmark. It was made on Opus 5.5 and has not been measured
% against the reference reading of folder 115 (made on Opus 5). The
% literature named in the footnotes (SGA 3, Serre 1968, Thomason 1987, EGA)
% is cited from memory and was not checked against the sources.
%
% Scope. The odd pages (1-23) are his notes; the even pages (2-24) are their
% versos, typed leaves of EGA IV §§ 18.6-18.9, published text, described
% only as the transcription's notes describe them (one section, no
% restatement). The hand of the ink corrections on the typed leaves is not
% established; said once.
%
% Spine. Equivariant modules, from a Dedekind ring to a scheme with an
% ample linearised line bundle:
%   p. 1        Lemma 1: saturation in a torsion-free module over a Dedekind
%               ring; a) <=> d) (filtered union of saturated f.g.
%               submodules); e) (linear forms separate) sufficient; necessary
%               for countably generated M. b), c) largely illegible.
%   p. 3        Lemma 2: every f.g. M in B lies in a f.g. saturated N with
%               Delta N in N (x) N.
%   p. 5        two corollaries: faithful linear representation over a
%               Dedekind ring; over a field, affine groups are strict
%               projective limits of affine groups of finite type.
%   pp. 7, 9    B = union of B_i (B/B_i flat, Delta B_i in B_i (x) B_i, B_i
%               f.g. projective): f.g. G-modules = union of the categories of
%               f.g. modules over U_i = dual of B_i; every f.g. G-module is a
%               quotient of a projective one.
%   pp. 9, 11   the same over a scheme S' with the resolution property.
%   pp. 13, 15  general statement 1)-4) for O_X-G-modules (read here with G
%               acting trivially on X).
%   pp. 17, 19  restatement with G acting on X and a G-linearised ample L;
%   pp. 19-23   proof of 1) and 2) through the graded module of sections.
%
% Notation fixed for the whole folder. Delta f(g,h) = f(gh) (comultiplication);
% a G-module is a comodule delta : E -> B (x) E, as on p. 7; U_i = \check B_i
% the dual of B_i; E_i the i-admissible part of E; N~ the saturation (pp.
% 1-5), M~ the sheaf of a graded module (pp. 19-23) — two different tildes,
% kept apart in the text; Gamma^.(E) = (+)_n f_*(E(n)) (written Gamma_* then
% Gamma^. on p. 19); A^. = Gamma^.(O_X).
%
% Substantive departures from the page, with pages.
%  - p. 3: « N of finite type » is supplied in Lemma 2 (without it N = B
%    works); the core of the proof is illegible, and the local-finiteness step
%    is supplied and attributed.
%  - p. 3: with Delta f(g,h) = f(gh), stability under G x e corresponds to
%    Delta N in B (x) N, not N (x) B; the page's « resp. » is swapped;
%    harmless, only the conjunction is used.
%  - p. 3: the tildes missing on the last three N of the final formula are
%    restored.
%  - pp. 13-15: no action of G on X is stated; read with trivial action (the
%    construction F (x) \check B_i of 3) needs it). The explicit formula
%    F' = sum_i U_i (F cap E_i) is this reading's.
%  - pp. 9, 15: « S (or X) noetherian regular » needs separated; X
%    quasi-compact supplied where a single index i is needed.
%  - p. 19: « il suffit de prouver 1) et 3) »: the proof of pp. 19-23 proves
%    1) and 2); 3) is never proved; a proof is sketched and marked ours.
%  - p. 23: « sous-module G-invariant de Gamma^n(F) » read as the smallest
%    G-stable submodule of Gamma^n(E) containing the phi_i.
%
% Announced and never established in the folder: conditions b), c) of Lemma
% 1 (illegible); the claim of p. 9 that Lemma 1's conditions on B mean that
% the irreducible components of G have non-empty fibres; the cancelled
% Proposition of p. 13; part 3) of p. 19; anything about K-groups — despite
% the title « KG (X) », no K-group, no Riemann-Roch appears.
\documentclass[11pt,a4paper]{article}
\input{../preamble/grothendieck}

\folder{90}
\pages{1}{24}
\dating{[à partir de 1964]}
\watermark{Édition de démonstration\\interprétation personnelle de l'œuvre}
\foldertitle{KG (X) : notes manuscrites (s.d.) — lecture modernisée du dossier entier}

\begin{document}

\section*{Résumé}

\begin{resume}
Un groupe agit sur un objet ; on voudrait ramener cette action à des
matrices, c'est-à-dire à une action sur des espaces de dimension finie. Sur
un corps, c'est une vieille idée : toute action « raisonnable » d'un groupe
algébrique sur un espace vectoriel est réunion de petites actions de
dimension finie, comme une fonction sur un groupe fini se décompose suivant
ses translatées. Ces douze pages manuscrites demandent ce qui reste de ce
principe quand le corps est remplacé par un anneau comme $\mathbf{Z}$, puis
par une base géométrique quelconque, et quand le groupe agit non plus sur un
module mais sur des faisceaux au-dessus d'un espace $X$ sur lequel il agit
lui-même.

La difficulté propre à $\mathbf{Z}$ est la torsion. Dans un espace vectoriel,
un sous-espace de dimension finie est fermé sur lui-même ; dans un groupe
abélien sans torsion, le sous-groupe engendré par un élément $x$ n'est pas
« saturé » : il faut lui adjoindre tous les $y$ dont un multiple $ny$ lui
appartient, et ce saturé peut cesser d'être de type fini (pensez à
$\mathbf{Z} \subset \mathbf{Q}$). Le premier lemme dit exactement quand ce
phénomène ne se produit pas, et le second s'en sert pour trouver, dans
l'algèbre des fonctions d'un groupe, des morceaux de type fini saturés et
stables par les translations. Deux corollaires en sortent : un groupe de
type fini sur un anneau de Dedekind a une représentation linéaire fidèle,
et sur un corps tout groupe affine est limite de groupes de type fini.

L'idée qui organise la suite est un changement de point de vue : une action
d'un groupe $G$ sur un module équivaut à une action d'une \emph{algèbre}, le
dual d'un morceau de l'algèbre des fonctions de $G$ — de même que, pour un
groupe fini, les représentations sont les modules sur l'algèbre du groupe.
Dès lors, tout ce qu'on sait sur les modules (être quotient d'un module
libre, par exemple) se transporte aux modules munis d'une action de $G$.
Grothendieck l'énonce sur une base $S$ quelconque, puis pour un schéma $X$
muni d'un fibré en droites ample compatible avec l'action : là, il ramène
les faisceaux sur $X$ à des modules gradués sur $S$, où l'action ne bouge
plus la base, et la question redevient algébrique. Le but est visible : tout
faisceau cohérent muni d'une action de $G$ est quotient d'un fibré vectoriel
muni d'une action de $G$, ce qui est la condition pour qu'une théorie $K$
équivariante — le $K_G(X)$ du titre — se calcule avec des fibrés. Mais le
dossier ne parle jamais de groupes $K$ : il pose les fondations et s'arrête
au milieu d'une démonstration, avec en marge « rédiger mieux ce début ».

Les noms d'aujourd'hui : comodules, sous-comodules et finitude locale ;
faisceaux $G$-équivariants et fibrés $G$-linéarisés ; propriété de
résolution équivariante, étudiée sous ce nom par Thomason.

\keywords{Dedekind domain, saturated submodule, Mittag-Leffler module, comodule, local finiteness of comodules, faithful linear representation, pro-algebraic group, dual algebra of a coalgebra, equivariant quasi-coherent sheaf, G-linearized ample line bundle, equivariant resolution property, graded module of sections}
\end{resume}

\pagerange{1}{24}
\section*{Le fil du dossier, et les conventions}

Le dossier a vingt-quatre pages. Les pages impaires portent les notes, sans
pagination de sa main ; les pages paires en sont les versos : des feuillets
dactylographiés d'EGA IV, §§ 18.6 à 18.9, sans rapport avec les notes, dont
on dit un mot à la fin. L'ordre de lecture est celui des pages impaires.

\begin{itemize}
  \item \textbf{Page 1 : Lemme 1.} Saturation dans un module sans torsion
    sur un anneau de Dedekind ; quand les saturés des sous-modules de type
    fini restent de type fini.
  \item \textbf{Page 3 : Lemme 2.} Dans l'algèbre $B$ d'un groupe affine,
    tout sous-module de type fini est contenu dans un sous-module de type
    fini, saturé, stable par la comultiplication.
  \item \textbf{Page 5 : deux corollaires.} Représentation linéaire fidèle
    sur un anneau de Dedekind ; groupes affines sur un corps comme limites
    projectives strictes de groupes de type fini.
  \item \textbf{Pages 7 et 9 : les algèbres $U_i$.} Si $B$ est réunion de
    sous-cogèbres $B_i$ projectives de type fini, les $G$-modules de type
    fini sont les modules de type fini sur les duaux $U_i$ ; tout $G$-module
    de type fini est quotient d'un $G$-module projectif de type fini.
  \item \textbf{Pages 9 et 11 : sur un schéma $S'$.} Même conclusion, avec
    « localement libre », sur un schéma où les faisceaux de type fini sont
    quotients de localement libres.
  \item \textbf{Pages 13 et 15 : énoncé général 1) à 4)} pour un
    $\mathcal{O}_X$-$G$-module quasi-cohérent ; lu ici avec $G$ agissant
    trivialement sur $X$.
  \item \textbf{Pages 17 et 19 : reprise} avec $G$ agissant sur $X$ et un
    faisceau inversible ample $G$-linéarisé ; \textbf{pages 19 à 23}, la
    démonstration de 1) et 2) par le module gradué des sections, qui
    s'arrête là.
\end{itemize}

La progression est cohérente : chaque étape réduit la suivante à la
précédente. Le Lemme 1 sert le Lemme 2, qui fournit les hypothèses des pages
7 et 13 ; et la démonstration des pages 19 à 23 ramène l'action sur $X$ à
une action sur la base $S$, où l'énoncé des pages 13 et 15 s'applique.

Le titre de l'inventaire, « KG (X) », désigne le groupe de Grothendieck des
faisceaux $G$-équivariants sur $X$. Aucune page ne le nomme : le dossier en
est le préalable — l'existence de résolutions équivariantes —, non la
théorie\footnote{Le classement du dossier dans le groupe « Géométrie et
topologie combinatoire » (69 à 102) ne tient pas à son contenu. Il ne
rencontre pas les dossiers voisins déjà modernisés — 81 (cartes et
découpages), 88 (cartes standards), 89 (ensembles 3-3) —, qui sont
combinatoires ; il est de la géométrie algébrique, comme le 94 selon sa
transcription.}.

Sur la date : les notes sont écrites au dos de feuillets dactylographiés
d'EGA IV, § 18. Si, comme d'ordinaire, le papier dactylographié a servi de
brouillon, elles sont postérieures à la frappe de ces feuillets (la partie
d'EGA IV qui les contient a paru en 1967) ; les pages elles-mêmes ne le
disent pas, et la datation « à partir de 1964 » est celle de l'inventaire.

\textbf{Conventions, valables pour tout le dossier.} $A$ est un anneau
commutatif, $G$ un $A$-groupe affine d'algèbre $B$, de comultiplication
$\Delta \colon B \to B \otimes_A B$, $\Delta f(g,h) = f(gh)$, et de coünité
$\varepsilon$. Un $G$-module (un « $A$-$G$-module ») est un $A$-module $E$
muni d'une coaction $\delta \colon E \to B \otimes E$, comme à la page 7 ;
c'est un comodule sur $B$. Pour une sous-cogèbre $B_i$ projective de type
fini, $U_i = \check{B}_i = \mathrm{Hom}_A(B_i, A)$ est son algèbre duale. Le
tilde a deux sens dans le dossier, et on les garde séparés : aux pages 1 à 5,
$\widetilde{N}$ est le saturé d'un sous-module ; aux pages 19 à 23,
$\widetilde{M}$ est le faisceau sur $X$ associé à un module gradué $M$. Sur
un schéma, $\underline{\Gamma}^{\cdot}(E) = \bigoplus_n f_{*}(E(n))$ et
$A^{\cdot} = \underline{\Gamma}^{\cdot}(\mathcal{O}_X)$.

\pagerange{1}{1}
\section*{Saturation sur un anneau de Dedekind (page 1)}

Soit $A$ un anneau de Dedekind et $M$ un $A$-module plat, c'est-à-dire sans
torsion. Pour un sous-module $N \subset M$, soit $\widetilde{N}$ l'image
inverse dans $M$ de la torsion de $M/N$ : c'est le \emph{saturé} de $N$, le
plus petit sous-module $N' \supset N$ tel que $M/N'$ soit sans torsion, et
$\widetilde{N}/N$ est de torsion\footnote{La ligne qui le dit commence par
« D.--- », lecture incertaine (définition ? démonstration ?).}. On dit $N$
\emph{saturé} si $\widetilde{N} = N$.

\textbf{Lemme 1.} Les conditions suivantes sont équivalentes :
\begin{itemize}
  \item[a)] pour tout sous-module de type fini $N$ de $M$, $\widetilde{N}$
    est de type fini (de façon équivalente, $\widetilde{N}/N$ est de
    longueur finie\footnote{« longueur » est une lecture incertaine. Elle
    donne l'énoncé juste : $\widetilde{N}/N$ est de torsion, et un module de
    torsion de type fini sur un anneau de Dedekind est de longueur finie.}) ;
  \item[d)] $M$ est réunion filtrante de sous-modules $M_i$ de type fini et
    saturés ; pour $M_i \subset M_j$, $M_i$ est alors facteur direct de
    $M_j$.
\end{itemize}
Elles sont satisfaites si
\begin{itemize}
  \item[e)] les formes linéaires sur $M$ séparent ses éléments,
    c'est-à-dire s'il existe une injection $M \hookrightarrow A^{I}$ ; c'est
    le cas en particulier si $M$ est projectif ;
\end{itemize}
et, si $M$ admet un ensemble dénombrable de générateurs, e) est aussi
nécessaire — et équivaut alors à « $M$ projectif »\footnote{La page annonce
quatre conditions équivalentes a) à d) et deux autres, b) et c), en
dépendent. Elles sont presque entièrement illisibles et refaites plusieurs
fois. Ce qu'on en lit : b) fait intervenir, pour un idéal maximal
$\mathfrak{p}$, l'anneau $A(\mathfrak{p}) = \bigcup_n \mathfrak{p}^{-n} =
\Gamma(S - \{\mathfrak{p}\}, \mathcal{O}_S)$, $S = \mathrm{Spec}\,A$, dans un
passage biffé ; c) porte sur les localisés $M_{\mathfrak{p}}$ et le corps
des fractions $K$ ; les deux sont précédés de « si $A$ n'est pas un corps »
(sur un corps, toutes les conditions sont trivialement vraies). On ne les
reconstitue pas. Une remarque, qui est de cette lecture, montre qu'une
condition purement locale ne peut être que nécessaire : pour $A =
\mathbf{Z}$, le sous-groupe $M$ de $\mathbf{Q}$ engendré par les $1/p$,
$p$ premier, a tous ses localisés $M_{(p)} = p^{-1}\mathbf{Z}_{(p)}$ de
type fini et ne contient aucun $\mathbf{Z}[1/p]$, et pourtant le saturé de
$\mathbf{Z}$ dans $M$ est $M$ entier, qui n'est pas de type fini. Le texte
de « d) » (« de type fini et "saturés", i.e.\ \ldots\ facteur direct dans
$M_j$ ») est ajouté en interligne, et la fin, « les applications de
transition injectives », est d'une lecture incertaine. « Moins fine », pour
« $M$ projectif », est aussi incertain ; la page le donne comme condition
plus forte que e), ce qui est juste.}.

\emph{Démonstration} (la page n'en donne pas ; celle-ci est de cette
lecture). a)~$\Rightarrow$~d) : les $\widetilde{N}$, $N$ de type fini,
forment un système filtrant de sous-modules saturés de type fini dont la
réunion est $M$. Si $M_i \subset M_j$ sont saturés, $M_j/M_i \subset M/M_i$
est sans torsion de type fini, donc projectif, et $M_i$ est facteur direct
de $M_j$. d)~$\Rightarrow$~a) : un $N$ de type fini est contenu dans un
$M_i$ ; $M/M_i$, réunion filtrante des $M_j/M_i$ sans torsion, est sans
torsion, donc $\widetilde{N} \subset M_i$, qui est noethérien.
e)~$\Rightarrow$~a) : si $N$ est de type fini, $N \otimes K$ est de
dimension finie, donc un nombre fini de coordonnées $A^{I} \to A^{k}$ sont
injectives sur $N \otimes K \supset \widetilde{N}$, et $\widetilde{N}$ est
un sous-module de $A^{k}$. Enfin, si $M$ est dénombrablement engendré et
satisfait d), on extrait une suite croissante $M_1 \subset M_2 \subset
\cdots$ de saturés de type fini de réunion $M$ ; chaque $M_{n+1}/M_n$ est
projectif, donc $M \simeq \bigoplus_n M_{n+1}/M_n$ est projectif et se
plonge dans un $A^{I}$.

On retiendra que, sous d), $M$ est limite inductive de modules projectifs
de type fini à transitions scindées : un module plat et de Mittag-Leffler,
dirait-on aujourd'hui\footnote{Le nom est moderne (Raynaud et Gruson, 1971)
et n'est pas sur la page ; il n'est employé ici que parce qu'il désigne
exactement la propriété dont la suite du dossier se sert : la
formation des sous-comodules commute aux intersections.}.

\pagerange{3}{3}
\section*{Sous-modules stables de l'algèbre d'un groupe (page 3)}

\textbf{Lemme 2.} Soit $A$ un anneau de Dedekind et $G$ un $A$-groupe
affine d'algèbre $B$ plate, satisfaisant aux conditions du Lemme 1. Pour tout
sous-module de type fini $M$ de $B$, il existe un sous-module $N \supset M$
de $B$, \emph{de type fini}\footnote{« de type fini » n'est pas dans
l'énoncé de la page ; sans lui, $N = B$ convient et le lemme est vide. La
démonstration le vise (« $\widetilde{N}$ \ldots\ de t.f. »).}, tel que
\begin{itemize}
  \item[a)] $N$ est saturé — ce qui entraîne que $N \otimes_A N \to B
    \otimes_A B$ est injectif ;
  \item[b)] $\Delta N \subset N \otimes N$.
\end{itemize}
Un tel $N$, sans torsion et de type fini, est projectif ; b) en fait une
sous-cogèbre de $B$.

\emph{Démonstration.} On fait opérer $G \times G$ sur $B$ par
\[
  \bigl((g,h)f\bigr)(s) = f(g^{-1} s h) .
\]
Un sous-module $N$ à quotient $B/N$ plat est stable par $G \times e$ (resp.\
$e \times G$) si et seulement si $\Delta N \subset B \otimes N$ (resp.\
$\Delta N \subset N \otimes B$)\footnote{La page écrit la correspondance
dans l'autre sens ($G \times e$ avec $N \otimes B$) ; c'est affaire de
convention sur $\Delta$, et seule la conjonction des deux conditions sert.}.
L'argument a trois temps, dont le premier est presque illisible sur la
page\footnote{Le milieu de la démonstration, une dizaine de lignes, est
réduit à quelques mots lus (« le passage de $N$ à $\widetilde{N}$ », « la
condition d'être invariant par $G \times G$ », « de t.f. ») ; l'ordre
proposé ici les respecte, mais l'énoncé de finitude locale qui ouvre
l'argument est restitué par cette lecture.} :
\begin{enumerate}
  \item $M$ est contenu dans un sous-module de type fini $N_0$ stable par
    $G \times G$ : c'est la finitude locale des comodules sur un anneau de
    Dedekind quand $B$ est plat\footnote{On trouve cet énoncé chez Serre,
    « Groupes de Grothendieck des schémas en groupes réductifs déployés »
    (Publ.\ IHÉS 34, 1968), § 1 ; citation de mémoire, non vérifiée, et sans
    rien dire de l'ordre des deux textes, que les pages ne permettent pas de
    fixer.} ;
  \item le saturé $N = \widetilde{N_0}$ est de type fini par le Lemme 1 a),
    et il reste stable : si $x \in N$ et $ax \in N_0$, $a \neq 0$, alors
    $a\,\Delta x \in N \otimes B$, et $(B \otimes B)/(N \otimes B) = (B/N)
    \otimes B$ est sans torsion, donc $\Delta x \in N \otimes B$ ; de même
    de l'autre côté ;
  \item comme $B/N$ est plat,
    \[
      \Delta N \subset (N \otimes B) \cap (B \otimes N) = N \otimes N
    \]
    (un élément de $N \otimes B$ nul dans $B \otimes (B/N)$ est nul dans
    $N \otimes (B/N)$, qui s'y injecte).
\end{enumerate}
C'est la formule finale de la page\footnote{Sur la page, les tildes sont
portés sur les deux premiers $N$ de la formule et manquent aux trois
derniers ; ils sont rétablis ici, puisque l'égalité n'est vraie que pour
le saturé. Une insertion marginale, en partie illisible, se termine par
« $N = \widetilde{N}$ ».}.

\pagerange{5}{5}
\section*{Deux corollaires (page 5)}

\textbf{Corollaire (représentation fidèle).} Soient $A$ et $G$ comme au
Lemme 2, et $G$ de plus de type fini sur $A$. Alors $G$ admet une
représentation linéaire fidèle sur un $A$-module projectif de type
fini\footnote{Les deux corollaires de la page 5, et celui de la page 9,
portent un titre souligné lu « Cor.\ Main » sous réserve ; ce peut être
« Corollaire ». « de plus » est ajouté en interligne.}.

En effet, $B$ est une $A$-algèbre de type fini ; appliquons le Lemme 2 à
un sous-module de type fini qui l'engendre. On obtient $N$ qui a) engendre
l'algèbre $B$, b) est saturé, donc projectif de type fini, c) vérifie
$\Delta N \subset N \otimes B$ : $N$ est un sous-comodule de $B$ pour la
représentation régulière. Un point $g$ de $G$ à valeurs dans une
$A$-algèbre $A'$ qui opère trivialement sur $N \otimes A'$ opère
trivialement sur l'algèbre $B \otimes A'$ qu'il engendre, donc est
l'unité\footnote{La phrase qui introduit la représentation (« Alors \ldots\
de $G$ ») est illisible ; on la lit comme la représentation de $G$ dans
$N$.}. On a même davantage : les coefficients matriciels de $N$ contiennent
$N$ (puisque $n = (\varepsilon \otimes \mathrm{id})\Delta n$), donc
engendrent $B$, et $G \to \mathrm{GL}(N)$ est une immersion
fermée\footnote{Remarque de cette lecture ; la page dit seulement
« fidèle ».}.

\textbf{Corollaire (groupes sur un corps).} Si $A$ est un corps $k$, tout
groupe affine sur $k$ est limite projective d'un système projectif strict
— à morphismes de transition fidèlement plats — de groupes affines
\emph{de type fini}\footnote{Le début de l'énoncé est illisible, et
« fid.\ plats » est d'une lecture incertaine. La démonstration n'est pas
écrite.}. C'est un résultat classique : $B$ est réunion filtrante de ses
sous-algèbres de Hopf de type fini (celle qu'engendrent un $N$ du Lemme 2 et
son image par l'antipode), et une algèbre de Hopf est fidèlement plate sur
toute sous-algèbre de Hopf.

\pagerange{7}{9}
\section*{Les modules sur les algèbres duales (pages 7 et 9)}

\textbf{Proposition.} Soit $A$ un anneau\footnote{Le début de l'énoncé est
illisible ; rien dans la suite n'exige plus qu'un anneau commutatif.}, $G$
un $A$-groupe affine d'algèbre $B$. On suppose que $B$ est réunion
filtrante de sous-modules $B_i$ tels que
\begin{itemize}
  \item[a)] $B/B_i$ est plat sur $A$, donc $B_i \otimes B_i \subset B
    \otimes B$ ;
  \item[b)] $\Delta B_i \subset B_i \otimes B_i$ ;
  \item[c)] les $B_i$ sont projectifs de type fini.
\end{itemize}
Par a), $B_i \otimes E \subset B \otimes E$ pour tout $A$-module $E$. Un
$G$-module $E$ est dit \emph{$i$-admissible} si sa coaction se factorise
en $\delta_i \colon E \to B_i \otimes E$. Alors :
\begin{enumerate}
  \item la catégorie des $G$-modules de type fini est la réunion, pour $i$
    croissant, de ses sous-catégories pleines $C_i(G)_f$ formées des
    $i$-admissibles ;
  \item par b) et c), $B_i$ est une cogèbre et $U_i = \check{B}_i$ une
    algèbre associative unitaire ; les $U_i$ forment un système projectif
    strict d'algèbres ;
  \item $C_i(G)_f$ est canoniquement équivalente à la catégorie des
    $U_i$-modules de type fini, et les foncteurs de transition $C_i \subset
    C_j$ correspondent à la restriction des scalaires le long de $U_j \to
    U_i$\footnote{« strict » et « de t.f. » sont ajoutés en interligne. Une
    première rédaction de 2 et 3 est barrée au bas de la page 7.}.
\end{enumerate}

Tout tient en peu de mots. Les images par $\delta$ d'un système fini de
générateurs sont des sommes finies, contenues dans un $B_i \otimes E$, d'où
1. La coünité se restreint à $B_i$, et la coassociativité aussi, puisque
$B_i \otimes B_i \otimes E \to B \otimes B \otimes E$ est injectif par a)
et c) ; une $B_i$-coaction sur $E$ est la même chose, par dualité des
projectifs de type fini, qu'une structure de $U_i$-module, $u \cdot e = (u
\otimes \mathrm{id})\,\delta_i(e)$. Pour $i \leq j$, $B_j/B_i$ est plat de
type fini, donc projectif, $B_i$ est facteur direct de $B_j$ et $U_j \to
U_i$ est surjectif : c'est le sens de « strict ».

\textbf{Corollaire.} Sous ces hypothèses, tout $G$-module $E$ de type fini
est quotient d'un $G$-module \emph{projectif} de type fini\footnote{« de
t.f.\ $E$ » est ajouté en interligne et « projectif », souligné, au-dessus
d'un mot biffé ; le mot « module » est d'une lecture incertaine.}.

Car $E$ est un $U_i$-module de type fini pour un $i$, donc quotient de
$U_i^{n}$, qui est un $U_i$-module projectif de type fini sur $A$, donc un
$G$-module dans $C_i(G)_f$. Autrement dit, si $A^{n} \to E$ est surjectif,
$A^{n} \otimes \check{B}_i \to E$, $x \otimes u \mapsto u \cdot x$, est
un épimorphisme de $G$-modules ; c'est la forme que prend l'énoncé 3) de la
page 15.

\pagerange{9}{11}
\section*{Sur un schéma de base (pages 9 et 11)}

\textbf{Proposition.} Soit $A$ un anneau de Dedekind, $S = \mathrm{Spec}\,A$,
$G$ un groupe affine plat sur $A$ d'algèbre $B$ satisfaisant aux conditions
du Lemme 1\footnote{La page ajoute entre parenthèses : « cela signifie que
pour les composantes irréductibles de $G$, \emph{leurs} fibres sont non
vides ». L'équivalence n'est pas démontrée et cette lecture ne l'a pas
vérifiée en général. Un exemple en montre le sens : sur $\mathbf{Z}_p$, le
groupe obtenu en ôtant du groupe constant $\mathbf{Z}/2$ le point non
neutre de la fibre spéciale a pour algèbre $\mathbf{Z}_p \times
\mathbf{Q}_p$ ; sa composante $\mathrm{Spec}\,\mathbf{Q}_p$ a une fibre
spéciale vide, et le saturé de $\mathbf{Z}_p \cdot (0,1)$ est $0 \times
\mathbf{Q}_p$, qui n'est pas de type fini.}. Soit $S'$ un schéma
quasi-compact sur $S$ sur lequel tout module quasi-cohérent de type fini
est quotient d'un module localement libre de type fini — par exemple $S'$
muni d'un faisceau inversible ample, ou noethérien régulier et
séparé\footnote{La page écrit « $S$ noeth.\ régulier », et un autre
exemple illisible. Pour un schéma noethérien régulier, la propriété vaut
s'il est séparé (Kleiman ; SGA 6, II, 2.2.7.1, de mémoire) ; la
séparation, comme la quasi-compacité qui permet de prendre un seul indice
$i$, est ajoutée ici. « schéma », « localement libre » et « p.\ ex. » sont
des lectures incertaines, et plusieurs mots de l'hypothèse sur $S'$ sont
illisibles.}. Alors tout $G_{S'}$-module quasi-cohérent de type fini est
quotient d'un $G_{S'}$-module quasi-cohérent de type fini qui est de plus
\emph{localement libre}.

La démonstration n'est pas écrite. Elle tient aux deux sections
précédentes : par le Lemme 2, $B$ est réunion filtrante de sous-cogèbres
$B_i$ saturées de type fini (on sature et on stabilise chaque sous-module
de type fini), qui satisfont a), b), c) de la page 7. Si $F \to E$ est un
épimorphisme d'un localement libre de type fini (sans action), $F \otimes
\check{B}_i \to E$ est un épimorphisme de $G$-modules dès que $E$ est
$i$-admissible, comme dans le corollaire de la page 9.

\pagerange{13}{15}
\section*{L'énoncé général (pages 13 et 15)}

La page 13 s'ouvre sur une Proposition barrée : pour $A$ de Dedekind et $G$
plat affine de type fini sur $A$, $N$ comme dans le Lemme 2 et la
représentation de $G$ dans $N$ induite par la représentation régulière
gauche, le morphisme $G \to \underline{\mathrm{Aut}}(N)$ aurait une
propriété dont le nom est perdu\footnote{La fin se lit sous réserve
« morphisme \ldots\ plat », peut-être « fidèle et plat ». Tout l'énoncé est
barré de six traits ; ce qu'on peut en dire de juste est l'immersion fermée
du premier corollaire de la page 5.}. Il l'abandonne pour un énoncé sur une
base quelconque.

Soit $S$ un schéma, $G$ un $S$-groupe affine d'algèbre $\mathcal{B}$
(quasi-cohérente sur $S$), réunion filtrante de sous-Modules
$\mathcal{B}_i$ tels que
\begin{itemize}
  \item[1)] $\mathcal{B}/\mathcal{B}_i$ est plat, donc $\mathcal{B}_i
    \otimes \mathcal{B}_i \subset \mathcal{B} \otimes \mathcal{B}$ ;
  \item[2)] $\Delta \mathcal{B}_i \subset \mathcal{B}_i \otimes
    \mathcal{B}_i$ ;
  \item[3)] les $\mathcal{B}_i$ sont localement libres de type fini.
\end{itemize}
Soit $X$ un $S$-schéma, et $E$ un $\mathcal{O}_X$-$G$-module
quasi-cohérent, $G$ agissant trivialement sur $X$\footnote{La page ne dit
rien d'une action de $G$ sur $X$. L'énoncé 3), où apparaît
$F \otimes \check{\mathcal{B}}_i$, n'a de sens tel quel que pour une
action triviale, et la reprise de la page 17 introduit l'action sur $X$ ;
on lit donc les pages 13 et 15 ainsi. C'est d'ailleurs le cas auquel la
démonstration des pages 19 à 23 ramène le cas général. Une hypothèse
« Supposons $X$ quasi-proj.\ \ldots » est biffée après 3).}. Alors :
\begin{itemize}
  \item[1)] tout sous-module quasi-cohérent $F$ de $E$ est contenu dans un
    plus petit sous-module quasi-cohérent $F'$ stable par $G$ ; sa
    formation commute à tout changement de base plat ; si $F$ est de type
    fini, $F'$ aussi ;
  \item[2)] si $X$ est quasi-compact et quasi-séparé, $E$ est réunion
    filtrante de sous-modules quasi-cohérents de type fini stables par $G$ ;
  \item[3)] si $X$ est quasi-compact, $F$ quasi-cohérent de type fini et
    $F \to E$ un morphisme dont l'image engendre $E$ comme $G$-module, il
    existe un épimorphisme de $G$-modules $F \otimes \check{\mathcal{B}}_i
    \to E$ — localement libre de type fini si $F$ l'est\footnote{Le début
    de 3) est illisible (« Soit $F \to E$ un \ldots ») ; « l'image
    engendre $E$ » est ce qu'on lit. La quasi-compacité de $X$, qui permet
    un seul $i$, est ajoutée. Un passage biffé entre 2) et 3) mentionne
    « $X$ quasi-compact et \ldots ».} ;
  \item[4)] si tout module quasi-cohérent de type fini sur $X$ est quotient
    d'un localement libre de type fini (par exemple si $X$ admet un faisceau
    inversible ample, ou s'il est noethérien régulier et séparé) et $X$
    quasi-compact, tout $\mathcal{O}_X$-$G$-module quasi-cohérent de type
    fini est quotient d'un $\mathcal{O}_X$-$G$-module localement libre de
    type fini\footnote{« quasi-proj. », « admet » et « régulier » sont des
    lectures incertaines ; « de type fini » pour le module de départ est
    restitué sous un mot illisible. Sur la séparation, voir la note de la
    page 9.}.
\end{itemize}

Aucune démonstration n'est écrite ; voici celle que les pages 7 et 9
contiennent en germe — la formule est de cette lecture. Soit $E_i \subset
E$ la partie $i$-admissible, noyau de $E \to (\mathcal{B}/\mathcal{B}_i)
\otimes E$ ; c'est un sous-$G$-module (avec une base duale locale $(b_k,
b_k^{*})$ de $\mathcal{B}_i$, $\delta e = \sum b_k \otimes b_k^{*} \cdot e$),
un $U_i$-module, et $E = \bigcup_i E_i$. Alors
\[
  F' = \sum_i U_i \cdot (F \cap E_i) .
\]
C'est un sous-$G$-module contenant $F$ (la coünité est dans $U_i$) ; et
tout sous-$G$-module $E''$ contenant $F$ est stable par chaque $U_i$, car la
platitude de $\mathcal{B}/\mathcal{B}_i$ donne $(\mathcal{B}_i \otimes E)
\cap (\mathcal{B} \otimes E'') = \mathcal{B}_i \otimes E''$. La formation de
$E_i$ et des images commute aux changements de base plats. Si $F$ est de
type fini, localement $F \subset E_{i_0}$ et $F' = U_{i_0} \cdot F$, image
de $\check{\mathcal{B}}_{i_0} \otimes F$, est de type fini — ce qui donne
3) ; 2) suit de 1), puisqu'un module quasi-cohérent sur un schéma
quasi-compact quasi-séparé est réunion de ses sous-modules de type fini ;
4) suit de 3).

\pagerange{17}{19}
\section*{Avec une action sur $X$ (pages 17 et 19)}

L'énoncé est repris, cette fois avec $G$ agissant sur $X$.

Soit $G$ un groupe plat et affine sur $S$, d'algèbre $\mathcal{B}$
satisfaisant 1), 2), 3) de la page 13 ; par 2) et 3), les
$\check{\mathcal{B}}_i$ sont des faisceaux d'algèbres associatives
unitaires, localement libres de type fini. Soit $f \colon X \to S$
quasi-compact et quasi-séparé\footnote{Hypothèse posée au cours de la
démonstration, page 19 (« $f$ qu.-cpt et (qu.-)sép. »).}, sur lequel $G$
opère, et $L$ un faisceau inversible ample relativement à $S$ sur lequel
$G$ opère de façon compatible à son action sur $X$ — un faisceau
inversible $G$-linéarisé, dit-on aujourd'hui\footnote{Le mot n'est pas sur
la page. La notion de linéarisation est dans la théorie géométrique des
invariants de Mumford (1965) ; on n'en tire rien sur ce que Grothendieck
connaissait en écrivant. Une première rédaction, « muni de $L$ ample
rel$/S$ », est biffée et remplacée par celle-ci.}. Soit $E$ un
$\mathcal{O}_X$-$G$-Module quasi-cohérent.
\begin{itemize}
  \item[1)] Pour tout sous-module quasi-cohérent $F$ de $E$, il existe un
    plus petit sous-module quasi-cohérent $F'$ de $E$ stable par $G$ et
    contenant $F$ ; sa formation commute à tout changement de base plat
    $S' \to S$.
  \item[2)] Si $F$ est de type fini, $F'$ l'est aussi\footnote{Écrit en
    interligne au-dessus d'un 2) barré : « si $S$ \ldots\ quasi-compact
    quasi-séparé, $E$ est limite de sous-modules quasi-cohérents de type
    fini stables par $G$ ». L'énoncé barré découle de 1) et 2), comme à la page 13.}.
  \item[3)] Si $S$ est quasi-compact, $E$ de type fini, et si tout module
    quasi-cohérent de type fini sur $S$ est quotient d'un localement libre
    de type fini, $E$ est quotient d'un $\mathcal{O}_X$-$G$-Module
    localement libre de type fini de la forme $f^{*}(H)(-n)$, $H$ un
    $G$-module localement libre de type fini sur $S$\footnote{La condition
    sur $S$ est écrite en interligne au-dessus d'une ligne biffée et
    commence par un mot lu « quasi-proj. » sous réserve. « $E$ de type
    fini » et « $S$ quasi-compact » sont ajoutés : sans eux la conclusion
    est fausse ou n'a pas d'indice uniforme. La version barrée de 3)
    affirmait sans condition l'existence d'un épimorphisme $F \to E$ avec
    $F$ localement libre $G$-équivariant ; au-dessus, ajouté puis biffé
    avec elle : « et est quotient d'un module loc.\ libre (sans op.\,!) »
    — c'est la condition qui manquait, et que la nouvelle version met sur
    $S$.}.
\end{itemize}

\pagerange{19}{23}
\section*{La démonstration par le module gradué des sections (pages 19 à 23)}

« Il suffit de prouver 1) et 3) » ; la démonstration écrite prouve 1) et
2), et s'arrête\footnote{La phrase annonce 1) et 3), peut-être selon une
numérotation antérieure à l'insertion du 2) de la page 17. 3) n'est
démontré nulle part dans le dossier ; une démonstration est esquissée à
la fin de cette section, et elle est de cette lecture.}.

\pagerange{19}{19}
\subsection*{Le dictionnaire gradué}

Posons
\[
  \underline{\Gamma}^{\cdot}(E) = \bigoplus_{n} f_{*}(E(n)), \qquad
  E(n) = E \otimes L^{\otimes n} .
\]
La page écrit d'abord $\underline{\Gamma}_{*}$ et un coproduit, puis
$\underline{\Gamma}^{\cdot}$ pour le même objet\footnote{« $\underline{\Gamma}_{*}(E)
= \coprod_n f_{*}(E(n))$ » ; on garde la seconde notation.}.
Comme $G$ est plat sur $S$ et $f$ quasi-compact et quasi-séparé,
$f_{*}$ commute au changement de base plat $G \to S$, et
$\underline{\Gamma}^{\cdot}(E)$ est un $\mathcal{O}_S$-Module gradué
quasi-cohérent sur lequel $G$ opère — \emph{trivialement sur la base} $S$.
De même $A^{\cdot} = \underline{\Gamma}^{\cdot}(\mathcal{O}_X)$ est une
$\mathcal{O}_S$-Algèbre graduée sur laquelle $G$ opère, et $E \mapsto
\underline{\Gamma}^{\cdot}(E)$ est un foncteur exact à gauche à valeurs dans
les $A^{\cdot}$-$G$-modules gradués. En sens inverse, $M \mapsto
\widetilde{M}$ est exact, adjoint à gauche de $\underline{\Gamma}^{\cdot}$,
et comme $L$ est ample, la coünité $\widetilde{\underline{\Gamma}^{\cdot}(E)}
\to E$ est un isomorphisme\footnote{EGA II, §§ 3 et 4 (référence de
mémoire) : pour $L$ ample relativement à $S$, $X$ s'envoie par une immersion
ouverte dans $\mathrm{Proj}(A^{\cdot})$ et $E \simeq
\widetilde{\underline{\Gamma}^{\cdot}(E)}$. « (ils sont adjoints \ldots) »
est ajouté en interligne, et « isom. » est d'une lecture incertaine.} ;
$\underline{\Gamma}^{\cdot}$ est donc pleinement fidèle. Tout se passe
désormais sur $S$, où l'action de $G$ ne déplace pas la base et où l'énoncé
des pages 13 et 15 s'applique.

\pagerange{21}{23}
\subsection*{Le plus petit sous-module stable (pages 21 et 23)}

Un sous-module quasi-cohérent $F \subset E$ correspond au sous-$A^{\cdot}$-module
$\underline{\Gamma}^{\cdot}(F)$ de $\underline{\Gamma}^{\cdot}(E)$, et si
$F \subset F_1 \subset E$ avec $F_1$ stable par $G$, alors
$\underline{\Gamma}^{\cdot}(F) \subset \underline{\Gamma}^{\cdot}(F_1)
\subset \underline{\Gamma}^{\cdot}(E)$ avec $\underline{\Gamma}^{\cdot}(F_1)$
stable par $G$\footnote{La page note $F'$ à la fois un sous-module stable
quelconque et le plus petit, qu'elle cherche ; on écrit ici $F_1$ pour le
premier. Plusieurs mots du passage sont illisibles, dont un mot inséré
au-dessus de « sous- » ; « à priori » et « s'identifie » sont incertains.}.

Par l'énoncé de la page 13, appliqué sur $S$, il y a un plus petit
sous-module $N$ de $\underline{\Gamma}^{\cdot}(E)$ contenant
$\underline{\Gamma}^{\cdot}(F)$ et stable par $G$. Il est \emph{gradué} :
l'action de $G$ respecte la graduation, chaque $U_i$ opère degré par degré,
et la formule $N = \sum_i U_i \cdot (\underline{\Gamma}^{\cdot}(F) \cap E_i)$
le montre\footnote{La page l'établit entre crochets, avec un point
d'interrogation final : $N \cap \underline{\Gamma}^{i}(E)$ contient
$\underline{\Gamma}^{i}(F)$ et est stable, donc contient le plus petit
sous-module stable $M^{i}$ de $\underline{\Gamma}^{i}(E)$ contenant
$\underline{\Gamma}^{i}(F)$ ; donc $N \supset \sum M^{i}$, et « par
minimalité lui est égal » (lecture incertaine). L'argument est juste,
pourvu que l'intersection de deux sous-modules stables le soit, ce que la
formule assure.}. C'est un sous-$A^{\cdot}$-module : la multiplication
$A^{\cdot} \otimes \underline{\Gamma}^{\cdot}(E) \to
\underline{\Gamma}^{\cdot}(E)$ est $G$-équivariante et
$\underline{\Gamma}^{\cdot}(F)$ est un $A^{\cdot}$-module, de sorte que,
pour des points, $a \cdot g x = g\bigl((g^{-1}a) \cdot x\bigr)$ reste dans
$N$ (« par fonctorialité de la construction de $N$ », dit la page).

Alors $\widetilde{N}$ est un $\mathcal{O}_X$-$G$-module quasi-cohérent, et
pour tout $F_1$ stable contenant $F$,
\[
  F = \widetilde{\underline{\Gamma}^{\cdot}(F)} \subset \widetilde{N}
  \subset \widetilde{\underline{\Gamma}^{\cdot}(F_1)} = F_1 .
\]
C'est donc le faisceau cherché, $F' = \widetilde{N}$. Par construction, sa
formation commute à toute extension plate $S' \to S$ de la base : c'est 1).

\emph{Finitude} (c'est 2)). Si $F$ est de type fini, $F'$ l'est : on peut
supposer $S$ affine, la question étant locale sur $S$ et la construction
compatible aux localisations. Comme $L$ est ample, il existe $n$ tel que
$F(n)$ soit engendré par un nombre fini de sections
$\varphi_1, \ldots, \varphi_r \in \Gamma(S, \underline{\Gamma}^{n}(F))$.
Soit $N^{n}$ le plus petit sous-module de $\underline{\Gamma}^{n}(E)$ stable
par $G$ et contenant les $\varphi_i$ ; il est de type fini par l'énoncé
1) de la page 13. Le sous-Module de $E$ qu'il définit, image de
$f^{*}(N^{n})(-n) \to E$, est stable par $G$, contient $F$ (engendré par
les $\varphi_i$) et est contenu dans $F'$ ; c'est donc $F'$, qui est de
type fini\footnote{La page dit « le sous-module $G$-invariant de
$\underline{\Gamma}^{n}(F)$ engendré par les $\varphi_i$ » : il faut lire
le plus petit sous-module stable de $\underline{\Gamma}^{n}(E)$, car
$\underline{\Gamma}^{n}(F)$ n'est pas stable. Une première tentative,
encadrée et barrée, prenait le sous-module engendré par les $\varphi_i$
sans le stabiliser. « $n$ » et « $F(n)$ » sont des lectures incertaines.
La page s'achève sur « On vérifie quand même que c'est $F'$, ce qui montre
que $F'$ est de type fini\ldots\ldots ».}.

En marge : « Rédiger mieux ce début de la démonstration, pour pouvoir y
\ldots »\footnote{Les deux derniers mots sont illisibles.} Le dossier
s'arrête là.

\emph{Ce que 3) demandait} (esquisse de cette lecture, non sur la page).
Si $E$ est de type fini, $E(n)$ est engendré par ses sections pour $n$
grand ; par 2) appliqué sur $S$, elles sont contenues dans un sous-$G$-module
$N^{n} \subset \underline{\Gamma}^{n}(E)$ de type fini, et par la page 9 (ou
le 4) de la page 15, sur $S$) $N^{n}$ est quotient d'un $G$-module $H$
localement libre de type fini. Alors $f^{*}(H)(-n) \to f^{*}(N^{n})(-n)
\to E$ est l'épimorphisme annoncé. C'est, sous le nom de « résolution
équivariante », un théorème de Thomason\footnote{R.~W.~Thomason,
« Equivariant resolution, linearization, and Hilbert's fourteenth problem
over arbitrary base schemes », Adv.\ in Math.\ 65 (1987) ; citation de
mémoire, non vérifiée. L'énoncé de la page 5 sur la représentation fidèle
se lit aussi, pour les groupes plats sur un anneau de Dedekind, dans SGA 3,
exposé VI$_B$, § 11 (de mémoire). Les pages ne citent rien, et l'on ne peut
pas dire d'après elles lesquels de ces textes elles précèdent ou suivent.}.

\pagerange{2}{24}
\section*{Les versos dactylographiés (pages paires)}

Les pages paires sont des feuillets de la frappe d'EGA IV, chapitre 18,
lus retournés ; le texte est publié et n'est pas repris ici. La
transcription les décrit un à un, avec les corrections à l'encre dont la
main n'est pas établie\footnote{Une marge de la page 22 emploie le signe
« ½ » pour \emph{semi-}, que l'on rencontre sous sa plume ; cela rapproche
ces corrections de lui sans les lui attribuer.}. Pour s'y retrouver :
\begin{itemize}
  \item pages 4 et 2 : « EGA IV 18.9. (Rabiots) », le Théorème 18.9.7 et les
    Corollaires 18.9.8 et 18.9.9, puis la fin d'une démonstration de 18.9.7
    (« caténaire » corrigé en un mot lu « japonais ») ;
  \item pages 6 et 14, paginées « -3- » et « -4- » : une rédaction reprise
    du Corollaire 18.9.5 ;
  \item page 8 (IV-965) : le Corollaire 18.9.4, une première rédaction
    barrée de 18.9.5, avec en marge « "séparés" inutile ! », et le début de
    la Remarque 18.9.6 ;
  \item page 10 (IV-964) : la Proposition 18.8.15, le Théorème 18.9.1, les
    Corollaires 18.9.2 et 18.9.3 ;
  \item page 12 : les Corollaires 18.8.12 et 18.8.13, la Proposition
    18.8.14, le début de 18.8.15, renumérotés à l'encre ;
  \item page 16 (IV-962) : la fin de 18.8.7, la Remarque 18.8.10, le début
    de la Proposition 18.8.11, avec un diagramme tracé à la main ;
  \item page 18 (IV-953) : les Corollaires 18.6.10 et 18.6.11 et la
    Proposition 18.6.12 ;
  \item page 20 : la fin du Théorème 18.6.9 et les Lemmes 18.6.9.1 et
    18.6.9.2, avec en marge « précisions mélangées ! » (lecture incertaine) ;
  \item pages 24 (IV-950) et 22 (IV-951) : la fin de 18.6.6, 18.6.7, la
    Proposition 18.6.8 et le début de 18.6.9 ; en marge de la page 22, une
    critique : « la forme hybride ici ne mérite pas le nom de "propr.\
    universelle" ».
\end{itemize}

\end{document}
